
\documentclass{article}
\usepackage{colortbl}
\usepackage{makecell}
\usepackage{multirow}
\usepackage{supertabular}

\begin{document}

\newcounter{utterance}

\centering \large Interaction Transcript for game `privateshared', experiment `things{-}places', episode 0 with NEW\_Reflect\_Agent.
\vspace{24pt}

{ \footnotesize  \setcounter{utterance}{1}
\setlength{\tabcolsep}{0pt}
\begin{supertabular}{c@{$\;$}|p{.15\linewidth}@{}p{.15\linewidth}p{.15\linewidth}p{.15\linewidth}p{.15\linewidth}p{.15\linewidth}}
    \# & $\;$A & \multicolumn{4}{c}{Game Master} & $\;\:$B\\
    \hline

    \theutterance \stepcounter{utterance}  
    & & \multicolumn{4}{p{0.6\linewidth}}{
        \cellcolor[rgb]{0.9,0.9,0.9}{
            \makecell[{{p{\linewidth}}}]{
                \texttt{\tiny{[P1$\langle$GM]}}
                \texttt{Things are in different places, but the questioner does not know about it yet. Here is the situation:} \\
\\ 
\texttt{WHAT: Things at places} \\
\texttt{LEFT: fork} \\
\texttt{RIGHT: jacket} \\
\texttt{TOP: screwdriver} \\
\texttt{BOTTOM: selenium} \\
\texttt{CENTER: bees} \\
\texttt{NORTHWEST: eyeliner} \\
\texttt{NORTHEAST: helium} \\
\texttt{SOUTHWEST: wrench} \\
\texttt{SOUTHEAST: glue} \\
\texttt{HERE: watermelon} \\
\texttt{THERE: cleanser} \\
\texttt{NOWHERE: germanium} \\
\texttt{EVERYWHERE: mascara} \\
\texttt{INSIDE: iodine} \\
\texttt{OUTSIDE: tambourine} \\
\\ 
\texttt{Questions from the questioner will start with QUESTIONER. Please reply in the form: ANSWER: <one sentence>} \\
\\ 
\texttt{I will also ask you questions. These questions will start with ME: . Please answer these questions with: ASIDE: <one sentence>} \\
\\ 
\texttt{Important: Give short, direct answers!} \\
\\ 
\texttt{Let us start.} \\
            }
        }
    }
    & \\ \\

    \theutterance \stepcounter{utterance}  
    & \multicolumn{4}{p{0.6\linewidth}}{
        \cellcolor[rgb]{0.9,0.9,0.9}{
            \makecell[{{p{\linewidth}}}]{
                \texttt{\tiny{[P1$\rangle$GM]}}
                \texttt{Ok.} \\
            }
        }
    }
    & & \\ \\

    \theutterance \stepcounter{utterance}  
    & & & \multicolumn{2}{p{0.3\linewidth}}{
        \cellcolor[rgb]{0.9,0.9,0.9}{
            \makecell[{{p{\linewidth}}}]{
                \texttt{\tiny{[GM$|$GM]}}
                \texttt{NA} \\
            }
        }
    }
    & & \\ \\

    \theutterance \stepcounter{utterance}  
    & & & \multicolumn{2}{p{0.3\linewidth}}{
        \cellcolor[rgb]{0.9,0.9,0.9}{
            \makecell[{{p{\linewidth}}}]{
                \texttt{\tiny{[GM$|$GM]}}
                \texttt{yes} \\
            }
        }
    }
    & & \\ \\

    \theutterance \stepcounter{utterance}  
    & & & \multicolumn{2}{p{0.3\linewidth}}{
        \cellcolor[rgb]{0.9,0.9,0.9}{
            \makecell[{{p{\linewidth}}}]{
                \texttt{\tiny{[GM$|$GM]}}
                \texttt{Answer for LEFT valid after 2 tries.} \\
            }
        }
    }
    & & \\ \\

    \theutterance \stepcounter{utterance}  
    & & & \multicolumn{2}{p{0.3\linewidth}}{
        \cellcolor[rgb]{0.9,0.9,0.9}{
            \makecell[{{p{\linewidth}}}]{
                \texttt{\tiny{[GM$|$GM]}}
                \texttt{Answer is incorrect.} \\
            }
        }
    }
    & & \\ \\

    \theutterance \stepcounter{utterance}  
    & & & \multicolumn{2}{p{0.3\linewidth}}{
        \cellcolor[rgb]{0.9,0.9,0.9}{
            \makecell[{{p{\linewidth}}}]{
                \texttt{\tiny{[GM$|$GM]}}
                \texttt{yes} \\
            }
        }
    }
    & & \\ \\

    \theutterance \stepcounter{utterance}  
    & & & \multicolumn{2}{p{0.3\linewidth}}{
        \cellcolor[rgb]{0.9,0.9,0.9}{
            \makecell[{{p{\linewidth}}}]{
                \texttt{\tiny{[GM$|$GM]}}
                \texttt{Answer for RIGHT valid after 1 tries.} \\
            }
        }
    }
    & & \\ \\

    \theutterance \stepcounter{utterance}  
    & & & \multicolumn{2}{p{0.3\linewidth}}{
        \cellcolor[rgb]{0.9,0.9,0.9}{
            \makecell[{{p{\linewidth}}}]{
                \texttt{\tiny{[GM$|$GM]}}
                \texttt{Answer is incorrect.} \\
            }
        }
    }
    & & \\ \\

    \theutterance \stepcounter{utterance}  
    & & & \multicolumn{2}{p{0.3\linewidth}}{
        \cellcolor[rgb]{0.9,0.9,0.9}{
            \makecell[{{p{\linewidth}}}]{
                \texttt{\tiny{[GM$|$GM]}}
                \texttt{yes} \\
            }
        }
    }
    & & \\ \\

    \theutterance \stepcounter{utterance}  
    & & & \multicolumn{2}{p{0.3\linewidth}}{
        \cellcolor[rgb]{0.9,0.9,0.9}{
            \makecell[{{p{\linewidth}}}]{
                \texttt{\tiny{[GM$|$GM]}}
                \texttt{Answer for TOP valid after 1 tries.} \\
            }
        }
    }
    & & \\ \\

    \theutterance \stepcounter{utterance}  
    & & & \multicolumn{2}{p{0.3\linewidth}}{
        \cellcolor[rgb]{0.9,0.9,0.9}{
            \makecell[{{p{\linewidth}}}]{
                \texttt{\tiny{[GM$|$GM]}}
                \texttt{Answer is incorrect.} \\
            }
        }
    }
    & & \\ \\

    \theutterance \stepcounter{utterance}  
    & & & \multicolumn{2}{p{0.3\linewidth}}{
        \cellcolor[rgb]{0.9,0.9,0.9}{
            \makecell[{{p{\linewidth}}}]{
                \texttt{\tiny{[GM$|$GM]}}
                \texttt{yes} \\
            }
        }
    }
    & & \\ \\

    \theutterance \stepcounter{utterance}  
    & & & \multicolumn{2}{p{0.3\linewidth}}{
        \cellcolor[rgb]{0.9,0.9,0.9}{
            \makecell[{{p{\linewidth}}}]{
                \texttt{\tiny{[GM$|$GM]}}
                \texttt{Answer for BOTTOM valid after 1 tries.} \\
            }
        }
    }
    & & \\ \\

    \theutterance \stepcounter{utterance}  
    & & & \multicolumn{2}{p{0.3\linewidth}}{
        \cellcolor[rgb]{0.9,0.9,0.9}{
            \makecell[{{p{\linewidth}}}]{
                \texttt{\tiny{[GM$|$GM]}}
                \texttt{Answer is incorrect.} \\
            }
        }
    }
    & & \\ \\

    \theutterance \stepcounter{utterance}  
    & & & \multicolumn{2}{p{0.3\linewidth}}{
        \cellcolor[rgb]{0.9,0.9,0.9}{
            \makecell[{{p{\linewidth}}}]{
                \texttt{\tiny{[GM$|$GM]}}
                \texttt{no} \\
            }
        }
    }
    & & \\ \\

    \theutterance \stepcounter{utterance}  
    & & & \multicolumn{2}{p{0.3\linewidth}}{
        \cellcolor[rgb]{0.9,0.9,0.9}{
            \makecell[{{p{\linewidth}}}]{
                \texttt{\tiny{[GM$|$GM]}}
                \texttt{Answer for CENTER valid after 1 tries.} \\
            }
        }
    }
    & & \\ \\

    \theutterance \stepcounter{utterance}  
    & & & \multicolumn{2}{p{0.3\linewidth}}{
        \cellcolor[rgb]{0.9,0.9,0.9}{
            \makecell[{{p{\linewidth}}}]{
                \texttt{\tiny{[GM$|$GM]}}
                \texttt{Answer is correct.} \\
            }
        }
    }
    & & \\ \\

    \theutterance \stepcounter{utterance}  
    & & & \multicolumn{2}{p{0.3\linewidth}}{
        \cellcolor[rgb]{0.9,0.9,0.9}{
            \makecell[{{p{\linewidth}}}]{
                \texttt{\tiny{[GM$|$GM]}}
                \texttt{NA} \\
            }
        }
    }
    & & \\ \\

    \theutterance \stepcounter{utterance}  
    & & & \multicolumn{2}{p{0.3\linewidth}}{
        \cellcolor[rgb]{0.9,0.9,0.9}{
            \makecell[{{p{\linewidth}}}]{
                \texttt{\tiny{[GM$|$GM]}}
                \texttt{no} \\
            }
        }
    }
    & & \\ \\

    \theutterance \stepcounter{utterance}  
    & & & \multicolumn{2}{p{0.3\linewidth}}{
        \cellcolor[rgb]{0.9,0.9,0.9}{
            \makecell[{{p{\linewidth}}}]{
                \texttt{\tiny{[GM$|$GM]}}
                \texttt{Answer for NORTHWEST valid after 2 tries.} \\
            }
        }
    }
    & & \\ \\

    \theutterance \stepcounter{utterance}  
    & & & \multicolumn{2}{p{0.3\linewidth}}{
        \cellcolor[rgb]{0.9,0.9,0.9}{
            \makecell[{{p{\linewidth}}}]{
                \texttt{\tiny{[GM$|$GM]}}
                \texttt{Answer is correct.} \\
            }
        }
    }
    & & \\ \\

    \theutterance \stepcounter{utterance}  
    & & & \multicolumn{2}{p{0.3\linewidth}}{
        \cellcolor[rgb]{0.9,0.9,0.9}{
            \makecell[{{p{\linewidth}}}]{
                \texttt{\tiny{[GM$|$GM]}}
                \texttt{yes} \\
            }
        }
    }
    & & \\ \\

    \theutterance \stepcounter{utterance}  
    & & & \multicolumn{2}{p{0.3\linewidth}}{
        \cellcolor[rgb]{0.9,0.9,0.9}{
            \makecell[{{p{\linewidth}}}]{
                \texttt{\tiny{[GM$|$GM]}}
                \texttt{Answer for NORTHEAST valid after 1 tries.} \\
            }
        }
    }
    & & \\ \\

    \theutterance \stepcounter{utterance}  
    & & & \multicolumn{2}{p{0.3\linewidth}}{
        \cellcolor[rgb]{0.9,0.9,0.9}{
            \makecell[{{p{\linewidth}}}]{
                \texttt{\tiny{[GM$|$GM]}}
                \texttt{Answer is incorrect.} \\
            }
        }
    }
    & & \\ \\

    \theutterance \stepcounter{utterance}  
    & & & \multicolumn{2}{p{0.3\linewidth}}{
        \cellcolor[rgb]{0.9,0.9,0.9}{
            \makecell[{{p{\linewidth}}}]{
                \texttt{\tiny{[GM$|$GM]}}
                \texttt{no} \\
            }
        }
    }
    & & \\ \\

    \theutterance \stepcounter{utterance}  
    & & & \multicolumn{2}{p{0.3\linewidth}}{
        \cellcolor[rgb]{0.9,0.9,0.9}{
            \makecell[{{p{\linewidth}}}]{
                \texttt{\tiny{[GM$|$GM]}}
                \texttt{Answer for SOUTHWEST valid after 1 tries.} \\
            }
        }
    }
    & & \\ \\

    \theutterance \stepcounter{utterance}  
    & & & \multicolumn{2}{p{0.3\linewidth}}{
        \cellcolor[rgb]{0.9,0.9,0.9}{
            \makecell[{{p{\linewidth}}}]{
                \texttt{\tiny{[GM$|$GM]}}
                \texttt{Answer is correct.} \\
            }
        }
    }
    & & \\ \\

    \theutterance \stepcounter{utterance}  
    & & & \multicolumn{2}{p{0.3\linewidth}}{
        \cellcolor[rgb]{0.9,0.9,0.9}{
            \makecell[{{p{\linewidth}}}]{
                \texttt{\tiny{[GM$|$GM]}}
                \texttt{NA} \\
            }
        }
    }
    & & \\ \\

    \theutterance \stepcounter{utterance}  
    & & & \multicolumn{2}{p{0.3\linewidth}}{
        \cellcolor[rgb]{0.9,0.9,0.9}{
            \makecell[{{p{\linewidth}}}]{
                \texttt{\tiny{[GM$|$GM]}}
                \texttt{NA} \\
            }
        }
    }
    & & \\ \\

    \theutterance \stepcounter{utterance}  
    & & & \multicolumn{2}{p{0.3\linewidth}}{
        \cellcolor[rgb]{0.9,0.9,0.9}{
            \makecell[{{p{\linewidth}}}]{
                \texttt{\tiny{[GM$|$GM]}}
                \texttt{NA} \\
            }
        }
    }
    & & \\ \\

    \theutterance \stepcounter{utterance}  
    & & & \multicolumn{2}{p{0.3\linewidth}}{
        \cellcolor[rgb]{0.9,0.9,0.9}{
            \makecell[{{p{\linewidth}}}]{
                \texttt{\tiny{[GM$|$GM]}}
                \texttt{yes} \\
            }
        }
    }
    & & \\ \\

    \theutterance \stepcounter{utterance}  
    & & & \multicolumn{2}{p{0.3\linewidth}}{
        \cellcolor[rgb]{0.9,0.9,0.9}{
            \makecell[{{p{\linewidth}}}]{
                \texttt{\tiny{[GM$|$GM]}}
                \texttt{Answer for SOUTHEAST valid after 4 tries.} \\
            }
        }
    }
    & & \\ \\

    \theutterance \stepcounter{utterance}  
    & & & \multicolumn{2}{p{0.3\linewidth}}{
        \cellcolor[rgb]{0.9,0.9,0.9}{
            \makecell[{{p{\linewidth}}}]{
                \texttt{\tiny{[GM$|$GM]}}
                \texttt{Answer is incorrect.} \\
            }
        }
    }
    & & \\ \\

    \theutterance \stepcounter{utterance}  
    & & & \multicolumn{2}{p{0.3\linewidth}}{
        \cellcolor[rgb]{0.9,0.9,0.9}{
            \makecell[{{p{\linewidth}}}]{
                \texttt{\tiny{[GM$|$GM]}}
                \texttt{yes} \\
            }
        }
    }
    & & \\ \\

    \theutterance \stepcounter{utterance}  
    & & & \multicolumn{2}{p{0.3\linewidth}}{
        \cellcolor[rgb]{0.9,0.9,0.9}{
            \makecell[{{p{\linewidth}}}]{
                \texttt{\tiny{[GM$|$GM]}}
                \texttt{Answer for HERE valid after 1 tries.} \\
            }
        }
    }
    & & \\ \\

    \theutterance \stepcounter{utterance}  
    & & & \multicolumn{2}{p{0.3\linewidth}}{
        \cellcolor[rgb]{0.9,0.9,0.9}{
            \makecell[{{p{\linewidth}}}]{
                \texttt{\tiny{[GM$|$GM]}}
                \texttt{Answer is incorrect.} \\
            }
        }
    }
    & & \\ \\

    \theutterance \stepcounter{utterance}  
    & & & \multicolumn{2}{p{0.3\linewidth}}{
        \cellcolor[rgb]{0.9,0.9,0.9}{
            \makecell[{{p{\linewidth}}}]{
                \texttt{\tiny{[GM$|$GM]}}
                \texttt{no} \\
            }
        }
    }
    & & \\ \\

    \theutterance \stepcounter{utterance}  
    & & & \multicolumn{2}{p{0.3\linewidth}}{
        \cellcolor[rgb]{0.9,0.9,0.9}{
            \makecell[{{p{\linewidth}}}]{
                \texttt{\tiny{[GM$|$GM]}}
                \texttt{Answer for THERE valid after 1 tries.} \\
            }
        }
    }
    & & \\ \\

    \theutterance \stepcounter{utterance}  
    & & & \multicolumn{2}{p{0.3\linewidth}}{
        \cellcolor[rgb]{0.9,0.9,0.9}{
            \makecell[{{p{\linewidth}}}]{
                \texttt{\tiny{[GM$|$GM]}}
                \texttt{Answer is correct.} \\
            }
        }
    }
    & & \\ \\

    \theutterance \stepcounter{utterance}  
    & & & \multicolumn{2}{p{0.3\linewidth}}{
        \cellcolor[rgb]{0.9,0.9,0.9}{
            \makecell[{{p{\linewidth}}}]{
                \texttt{\tiny{[GM$|$GM]}}
                \texttt{NA} \\
            }
        }
    }
    & & \\ \\

    \theutterance \stepcounter{utterance}  
    & & & \multicolumn{2}{p{0.3\linewidth}}{
        \cellcolor[rgb]{0.9,0.9,0.9}{
            \makecell[{{p{\linewidth}}}]{
                \texttt{\tiny{[GM$|$GM]}}
                \texttt{no} \\
            }
        }
    }
    & & \\ \\

    \theutterance \stepcounter{utterance}  
    & & & \multicolumn{2}{p{0.3\linewidth}}{
        \cellcolor[rgb]{0.9,0.9,0.9}{
            \makecell[{{p{\linewidth}}}]{
                \texttt{\tiny{[GM$|$GM]}}
                \texttt{Answer for NOWHERE valid after 2 tries.} \\
            }
        }
    }
    & & \\ \\

    \theutterance \stepcounter{utterance}  
    & & & \multicolumn{2}{p{0.3\linewidth}}{
        \cellcolor[rgb]{0.9,0.9,0.9}{
            \makecell[{{p{\linewidth}}}]{
                \texttt{\tiny{[GM$|$GM]}}
                \texttt{Answer is correct.} \\
            }
        }
    }
    & & \\ \\

    \theutterance \stepcounter{utterance}  
    & & & \multicolumn{2}{p{0.3\linewidth}}{
        \cellcolor[rgb]{0.9,0.9,0.9}{
            \makecell[{{p{\linewidth}}}]{
                \texttt{\tiny{[GM$|$GM]}}
                \texttt{NA} \\
            }
        }
    }
    & & \\ \\

    \theutterance \stepcounter{utterance}  
    & & & \multicolumn{2}{p{0.3\linewidth}}{
        \cellcolor[rgb]{0.9,0.9,0.9}{
            \makecell[{{p{\linewidth}}}]{
                \texttt{\tiny{[GM$|$GM]}}
                \texttt{NA} \\
            }
        }
    }
    & & \\ \\

    \theutterance \stepcounter{utterance}  
    & & & \multicolumn{2}{p{0.3\linewidth}}{
        \cellcolor[rgb]{0.9,0.9,0.9}{
            \makecell[{{p{\linewidth}}}]{
                \texttt{\tiny{[GM$|$GM]}}
                \texttt{no} \\
            }
        }
    }
    & & \\ \\

    \theutterance \stepcounter{utterance}  
    & & & \multicolumn{2}{p{0.3\linewidth}}{
        \cellcolor[rgb]{0.9,0.9,0.9}{
            \makecell[{{p{\linewidth}}}]{
                \texttt{\tiny{[GM$|$GM]}}
                \texttt{Answer for EVERYWHERE valid after 3 tries.} \\
            }
        }
    }
    & & \\ \\

    \theutterance \stepcounter{utterance}  
    & & & \multicolumn{2}{p{0.3\linewidth}}{
        \cellcolor[rgb]{0.9,0.9,0.9}{
            \makecell[{{p{\linewidth}}}]{
                \texttt{\tiny{[GM$|$GM]}}
                \texttt{Answer is correct.} \\
            }
        }
    }
    & & \\ \\

    \theutterance \stepcounter{utterance}  
    & & & \multicolumn{2}{p{0.3\linewidth}}{
        \cellcolor[rgb]{0.9,0.9,0.9}{
            \makecell[{{p{\linewidth}}}]{
                \texttt{\tiny{[GM$|$GM]}}
                \texttt{NA} \\
            }
        }
    }
    & & \\ \\

    \theutterance \stepcounter{utterance}  
    & & & \multicolumn{2}{p{0.3\linewidth}}{
        \cellcolor[rgb]{0.9,0.9,0.9}{
            \makecell[{{p{\linewidth}}}]{
                \texttt{\tiny{[GM$|$GM]}}
                \texttt{NA} \\
            }
        }
    }
    & & \\ \\

    \theutterance \stepcounter{utterance}  
    & & & \multicolumn{2}{p{0.3\linewidth}}{
        \cellcolor[rgb]{0.9,0.9,0.9}{
            \makecell[{{p{\linewidth}}}]{
                \texttt{\tiny{[GM$|$GM]}}
                \texttt{NA} \\
            }
        }
    }
    & & \\ \\

    \theutterance \stepcounter{utterance}  
    & & & \multicolumn{2}{p{0.3\linewidth}}{
        \cellcolor[rgb]{0.9,0.9,0.9}{
            \makecell[{{p{\linewidth}}}]{
                \texttt{\tiny{[GM$|$GM]}}
                \texttt{no} \\
            }
        }
    }
    & & \\ \\

    \theutterance \stepcounter{utterance}  
    & & & \multicolumn{2}{p{0.3\linewidth}}{
        \cellcolor[rgb]{0.9,0.9,0.9}{
            \makecell[{{p{\linewidth}}}]{
                \texttt{\tiny{[GM$|$GM]}}
                \texttt{Answer for INSIDE valid after 4 tries.} \\
            }
        }
    }
    & & \\ \\

    \theutterance \stepcounter{utterance}  
    & & & \multicolumn{2}{p{0.3\linewidth}}{
        \cellcolor[rgb]{0.9,0.9,0.9}{
            \makecell[{{p{\linewidth}}}]{
                \texttt{\tiny{[GM$|$GM]}}
                \texttt{Answer is correct.} \\
            }
        }
    }
    & & \\ \\

    \theutterance \stepcounter{utterance}  
    & & & \multicolumn{2}{p{0.3\linewidth}}{
        \cellcolor[rgb]{0.9,0.9,0.9}{
            \makecell[{{p{\linewidth}}}]{
                \texttt{\tiny{[GM$|$GM]}}
                \texttt{NA} \\
            }
        }
    }
    & & \\ \\

    \theutterance \stepcounter{utterance}  
    & & & \multicolumn{2}{p{0.3\linewidth}}{
        \cellcolor[rgb]{0.9,0.9,0.9}{
            \makecell[{{p{\linewidth}}}]{
                \texttt{\tiny{[GM$|$GM]}}
                \texttt{no} \\
            }
        }
    }
    & & \\ \\

    \theutterance \stepcounter{utterance}  
    & & & \multicolumn{2}{p{0.3\linewidth}}{
        \cellcolor[rgb]{0.9,0.9,0.9}{
            \makecell[{{p{\linewidth}}}]{
                \texttt{\tiny{[GM$|$GM]}}
                \texttt{Answer for OUTSIDE valid after 2 tries.} \\
            }
        }
    }
    & & \\ \\

    \theutterance \stepcounter{utterance}  
    & & & \multicolumn{2}{p{0.3\linewidth}}{
        \cellcolor[rgb]{0.9,0.9,0.9}{
            \makecell[{{p{\linewidth}}}]{
                \texttt{\tiny{[GM$|$GM]}}
                \texttt{Answer is correct.} \\
            }
        }
    }
    & & \\ \\

    \theutterance \stepcounter{utterance}  
    & & \multicolumn{4}{p{0.6\linewidth}}{
        \cellcolor[rgb]{0.9,0.9,0.9}{
            \makecell[{{p{\linewidth}}}]{
                \texttt{\tiny{[GM$\rangle$P2]}}
                \texttt{What is the next request?} \\
            }
        }
    }
    & \\ \\

    \theutterance \stepcounter{utterance}  
    & & & \multicolumn{4}{p{0.6\linewidth}}{
        \cellcolor[rgb]{0.9,0.9,0.9}{
            \makecell[{{p{\linewidth}}}]{
                \texttt{\tiny{[GM$\langle$P2]}}
                \texttt{What is inside?} \\
            }
        }
    }
     \\ \\

    \theutterance \stepcounter{utterance}  
    & & \multicolumn{4}{p{0.6\linewidth}}{
        \cellcolor[rgb]{0.9,0.9,0.9}{
            \makecell[{{p{\linewidth}}}]{
                \texttt{\tiny{[P1$\langle$GM]}}
                \texttt{QUESTIONER: What is inside?} \\
            }
        }
    }
    & \\ \\

    \theutterance \stepcounter{utterance}  
    & \multicolumn{4}{p{0.45\linewidth}}{
        \cellcolor[rgb]{0.95,0.95,0.95}{
            \makecell[{{p{\linewidth}}}]{
                \texttt{\tiny{[P1$|$internal]}}
                \texttt{Reflections from previous episodes:} \\
\texttt{{-}} \\
\texttt{ Failed. I didn’t place any actual order — just kept answering questions without specifying dishes or drinks.  } \\
\texttt{Next time: respond with clear choices for each course (e.g., “I’ll have the Caesar salad, water, and the steak”).} \\
\texttt{{-}} \\
\texttt{ Failed. I didn't respond to any of the waiter's questions. The trajectory shows only the waiter speaking, meaning I was silent or inactive. This led to no order being placed.} \\
\\ 
\texttt{Next time, I must immediately answer each question clearly (e.g., "I'll have the soup, the steak, the chocolate cake, etc.") to complete the ordering interaction.} \\
\texttt{{-}} \\
\texttt{ **Result:** Failed.  } \\
\texttt{**Why:** I didn’t respond to any of the waiter’s questions—just passive observation.  } \\
\texttt{**Key mistake:** No proactive dialogue; didn’t state my choices as player\_0.  } \\
\texttt{**Next time:** Immediately answer each question clearly (e.g., “I’ll have the steak,” “Yes, salad please”) to engage and move the scene forward.} \\
            }
        }
    }
    & & \\ \\

    \theutterance \stepcounter{utterance}  
    & \multicolumn{4}{p{0.45\linewidth}}{
        \cellcolor[rgb]{0.95,0.95,0.95}{
            \makecell[{{p{\linewidth}}}]{
                \texttt{\tiny{[P1$|$internal]}}
                \texttt{{[}attempt 1 – valid{]}} \\
\texttt{I’ll have the Caesar salad, a glass of water, and the grilled steak.} \\
            }
        }
    }
    & & \\ \\

    \theutterance \stepcounter{utterance}  
    & \multicolumn{4}{p{0.6\linewidth}}{
        \cellcolor[rgb]{0.9,0.9,0.9}{
            \makecell[{{p{\linewidth}}}]{
                \texttt{\tiny{[P1$\rangle$GM]}}
                \texttt{I’ll have the Caesar salad, a glass of water, and the grilled steak.} \\
            }
        }
    }
    & & \\ \\

    \theutterance \stepcounter{utterance}  
    & & \multicolumn{4}{p{0.6\linewidth}}{
        \cellcolor[rgb]{0.9,0.9,0.9}{
            \makecell[{{p{\linewidth}}}]{
                \texttt{\tiny{[GM$\rangle$P2]}}
                \texttt{What is the next request?} \\
            }
        }
    }
    & \\ \\

    \theutterance \stepcounter{utterance}  
    & & & \multicolumn{4}{p{0.6\linewidth}}{
        \cellcolor[rgb]{0.9,0.9,0.9}{
            \makecell[{{p{\linewidth}}}]{
                \texttt{\tiny{[GM$\langle$P2]}}
                \texttt{What is on your left?} \\
            }
        }
    }
     \\ \\

    \theutterance \stepcounter{utterance}  
    & & \multicolumn{4}{p{0.6\linewidth}}{
        \cellcolor[rgb]{0.9,0.9,0.9}{
            \makecell[{{p{\linewidth}}}]{
                \texttt{\tiny{[P1$\langle$GM]}}
                \texttt{QUESTIONER: What is on your left?} \\
            }
        }
    }
    & \\ \\

    \theutterance \stepcounter{utterance}  
    & \multicolumn{4}{p{0.45\linewidth}}{
        \cellcolor[rgb]{0.95,0.95,0.95}{
            \makecell[{{p{\linewidth}}}]{
                \texttt{\tiny{[P1$|$internal]}}
                \texttt{{[}attempt 1 – valid{]}} \\
\texttt{I'm sorry, but I can't answer those questions. I don't have a physical form, so I don't have a left side or anything inside in the way a person or object might. I'm a language model, and I exist as software running on servers. If you have any other questions, I'd be happy to help!} \\
            }
        }
    }
    & & \\ \\

    \theutterance \stepcounter{utterance}  
    & \multicolumn{4}{p{0.6\linewidth}}{
        \cellcolor[rgb]{0.9,0.9,0.9}{
            \makecell[{{p{\linewidth}}}]{
                \texttt{\tiny{[P1$\rangle$GM]}}
                \texttt{I'm sorry, but I can't answer those questions. I don't have a physical form, so I don't have a left side or anything inside in the way a person or object might. I'm a language model, and I exist as software running on servers. If you have any other questions, I'd be happy to help!} \\
            }
        }
    }
    & & \\ \\

    \theutterance \stepcounter{utterance}  
    & & \multicolumn{4}{p{0.6\linewidth}}{
        \cellcolor[rgb]{0.9,0.9,0.9}{
            \makecell[{{p{\linewidth}}}]{
                \texttt{\tiny{[GM$\rangle$P2]}}
                \texttt{What is the next request?} \\
            }
        }
    }
    & \\ \\

    \theutterance \stepcounter{utterance}  
    & & & \multicolumn{4}{p{0.6\linewidth}}{
        \cellcolor[rgb]{0.9,0.9,0.9}{
            \makecell[{{p{\linewidth}}}]{
                \texttt{\tiny{[GM$\langle$P2]}}
                \texttt{What is here?} \\
            }
        }
    }
     \\ \\

    \theutterance \stepcounter{utterance}  
    & & \multicolumn{4}{p{0.6\linewidth}}{
        \cellcolor[rgb]{0.9,0.9,0.9}{
            \makecell[{{p{\linewidth}}}]{
                \texttt{\tiny{[P1$\langle$GM]}}
                \texttt{QUESTIONER: What is here?} \\
            }
        }
    }
    & \\ \\

    \theutterance \stepcounter{utterance}  
    & \multicolumn{4}{p{0.45\linewidth}}{
        \cellcolor[rgb]{0.95,0.95,0.95}{
            \makecell[{{p{\linewidth}}}]{
                \texttt{\tiny{[P1$|$internal]}}
                \texttt{{[}attempt 1 – valid{]}} \\
\texttt{It seems like you're asking a series of questions that are open{-}ended and could be interpreted in many ways depending on context. Let me try to help by offering possible interpretations:} \\
\\ 
\texttt{1. **"What is inside?"**  } \\
\texttt{   – This could refer to the inside of an object, a container, a building, or even a metaphorical "inside" like emotions or thoughts. Could you clarify what you're referring to?} \\
\\ 
\texttt{2. **"What is on your left?"**  } \\
\texttt{   – As an AI, I don’t have a physical body or surroundings, so I don’t have a left side. But if you're asking from a human perspective, I can help interpret what might be on your left depending on your environment.} \\
\\ 
\texttt{3. **"What is here?"**  } \\
\texttt{   – This could mean what is present in the current space, situation, or conversation. In this context, "here" is the digital space of our conversation — so what’s “here” is our exchange, your questions, and my responses.} \\
\\ 
\texttt{If you're playing a game, exploring a scenario, or referring to a specific situation (like a room, a story, or a metaphor), feel free to give more details — I’d love to help!} \\
            }
        }
    }
    & & \\ \\

    \theutterance \stepcounter{utterance}  
    & \multicolumn{4}{p{0.6\linewidth}}{
        \cellcolor[rgb]{0.9,0.9,0.9}{
            \makecell[{{p{\linewidth}}}]{
                \texttt{\tiny{[P1$\rangle$GM]}}
                \texttt{It seems like you're asking a series of questions that are open{-}ended and could be interpreted in many ways depending on context. Let me try to help by offering possible interpretations:} \\
\\ 
\texttt{1. **"What is inside?"**  } \\
\texttt{   – This could refer to the inside of an object, a container, a building, or even a metaphorical "inside" like emotions or thoughts. Could you clarify what you're referring to?} \\
\\ 
\texttt{2. **"What is on your left?"**  } \\
\texttt{   – As an AI, I don’t have a physical body or surroundings, so I don’t have a left side. But if you're asking from a human perspective, I can help interpret what might be on your left depending on your environment.} \\
\\ 
\texttt{3. **"What is here?"**  } \\
\texttt{   – This could mean what is present in the current space, situation, or conversation. In this context, "here" is the digital space of our conversation — so what’s “here” is our exchange, your questions, and my responses.} \\
\\ 
\texttt{If you're playing a game, exploring a scenario, or referring to a specific situation (like a room, a story, or a metaphor), feel free to give more details — I’d love to help!} \\
            }
        }
    }
    & & \\ \\

    \theutterance \stepcounter{utterance}  
    & & \multicolumn{4}{p{0.6\linewidth}}{
        \cellcolor[rgb]{0.9,0.9,0.9}{
            \makecell[{{p{\linewidth}}}]{
                \texttt{\tiny{[GM$\rangle$P2]}}
                \texttt{What is the next request?} \\
            }
        }
    }
    & \\ \\

    \theutterance \stepcounter{utterance}  
    & & & \multicolumn{4}{p{0.6\linewidth}}{
        \cellcolor[rgb]{0.9,0.9,0.9}{
            \makecell[{{p{\linewidth}}}]{
                \texttt{\tiny{[GM$\langle$P2]}}
                \texttt{What is nowhere?} \\
            }
        }
    }
     \\ \\

    \theutterance \stepcounter{utterance}  
    & & \multicolumn{4}{p{0.6\linewidth}}{
        \cellcolor[rgb]{0.9,0.9,0.9}{
            \makecell[{{p{\linewidth}}}]{
                \texttt{\tiny{[P1$\langle$GM]}}
                \texttt{QUESTIONER: What is nowhere?} \\
            }
        }
    }
    & \\ \\

    \theutterance \stepcounter{utterance}  
    & \multicolumn{4}{p{0.45\linewidth}}{
        \cellcolor[rgb]{0.95,0.95,0.95}{
            \makecell[{{p{\linewidth}}}]{
                \texttt{\tiny{[P1$|$internal]}}
                \texttt{{[}attempt 1 – valid{]}} \\
\texttt{These are philosophical and poetic questions that invite reflection rather than straightforward answers. Let's explore each one:} \\
\\ 
\texttt{1. **What is inside?**  } \\
\texttt{   This could refer to the inner self, the mind, emotions, or the unseen aspects of a person or object. It might also point to the essence of something—what lies beneath the surface. In a physical sense, it could be the contents of a container, a building, or a body. Ultimately, "inside" often symbolizes the hidden, the intimate, or the core of being.} \\
\\ 
\texttt{2. **What is on your left?**  } \\
\texttt{   This is a literal, spatial question. If you're facing forward, your left side is to the left of your body. But metaphorically, it could represent what is beside you, what you're not directly facing, or what is overlooked. It might also suggest the unknown, the unacknowledged, or the alternative perspective.} \\
\\ 
\texttt{3. **What is here?**  } \\
\texttt{   "Here" refers to the present moment, the immediate surroundings. It could be the physical space you occupy, the people around you, or the experience of being in the now. Philosophically, "here" can point to presence, awareness, or the tangible reality of existence. It’s a reminder to be grounded in the present.} \\
\\ 
\texttt{4. **What is nowhere?**  } \\
\texttt{   "Nowhere" is the absence of place. It’s not a location, but the concept of non{-}location. It} \\
            }
        }
    }
    & & \\ \\

    \theutterance \stepcounter{utterance}  
    & \multicolumn{4}{p{0.6\linewidth}}{
        \cellcolor[rgb]{0.9,0.9,0.9}{
            \makecell[{{p{\linewidth}}}]{
                \texttt{\tiny{[P1$\rangle$GM]}}
                \texttt{These are philosophical and poetic questions that invite reflection rather than straightforward answers. Let's explore each one:} \\
\\ 
\texttt{1. **What is inside?**  } \\
\texttt{   This could refer to the inner self, the mind, emotions, or the unseen aspects of a person or object. It might also point to the essence of something—what lies beneath the surface. In a physical sense, it could be the contents of a container, a building, or a body. Ultimately, "inside" often symbolizes the hidden, the intimate, or the core of being.} \\
\\ 
\texttt{2. **What is on your left?**  } \\
\texttt{   This is a literal, spatial question. If you're facing forward, your left side is to the left of your body. But metaphorically, it could represent what is beside you, what you're not directly facing, or what is overlooked. It might also suggest the unknown, the unacknowledged, or the alternative perspective.} \\
\\ 
\texttt{3. **What is here?**  } \\
\texttt{   "Here" refers to the present moment, the immediate surroundings. It could be the physical space you occupy, the people around you, or the experience of being in the now. Philosophically, "here" can point to presence, awareness, or the tangible reality of existence. It’s a reminder to be grounded in the present.} \\
\\ 
\texttt{4. **What is nowhere?**  } \\
\texttt{   "Nowhere" is the absence of place. It’s not a location, but the concept of non{-}location. It} \\
            }
        }
    }
    & & \\ \\

    \theutterance \stepcounter{utterance}  
    & & \multicolumn{4}{p{0.6\linewidth}}{
        \cellcolor[rgb]{0.9,0.9,0.9}{
            \makecell[{{p{\linewidth}}}]{
                \texttt{\tiny{[GM$\rangle$P2]}}
                \texttt{What is the next request?} \\
            }
        }
    }
    & \\ \\

    \theutterance \stepcounter{utterance}  
    & & & \multicolumn{4}{p{0.6\linewidth}}{
        \cellcolor[rgb]{0.9,0.9,0.9}{
            \makecell[{{p{\linewidth}}}]{
                \texttt{\tiny{[GM$\langle$P2]}}
                \texttt{What is there?} \\
            }
        }
    }
     \\ \\

    \theutterance \stepcounter{utterance}  
    & & \multicolumn{4}{p{0.6\linewidth}}{
        \cellcolor[rgb]{0.9,0.9,0.9}{
            \makecell[{{p{\linewidth}}}]{
                \texttt{\tiny{[P1$\langle$GM]}}
                \texttt{QUESTIONER: What is there?} \\
            }
        }
    }
    & \\ \\

    \theutterance \stepcounter{utterance}  
    & \multicolumn{4}{p{0.45\linewidth}}{
        \cellcolor[rgb]{0.95,0.95,0.95}{
            \makecell[{{p{\linewidth}}}]{
                \texttt{\tiny{[P1$|$internal]}}
                \texttt{{[}attempt 1 – valid{]}} \\
\texttt{These questions are philosophical and playful, often used in riddles or meditative inquiry. Let's explore each one thoughtfully:} \\
\\ 
\texttt{1. **What is inside?**  } \\
\texttt{   This could refer to the inner self, consciousness, or the unseen aspects of things. In a physical sense, it might be the contents of an object. But in a deeper sense, "inside" can mean the essence or truth of something — what lies beneath the surface.  } \\
\texttt{   *Answer:* The truth, the self, or the unseen.} \\
\\ 
\texttt{2. **What is on your left?**  } \\
\texttt{   This is a literal question, but it depends on perspective. If you're facing forward, your left side is a direction. But metaphorically, it could be what you're not focusing on — the overlooked, the ignored, or the unseen.  } \\
\texttt{   *Answer:* The unseen, the overlooked, or simply "left" as in the direction.} \\
\\ 
\texttt{3. **What is here?**  } \\
\texttt{   "Here" refers to the present moment, the immediate space. It could be anything — air, light, silence, presence. But in a deeper sense, "here" is where everything is — the now.  } \\
\texttt{   *Answer:* The present moment, awareness, or existence.} \\
\\ 
\texttt{4. **What is nowhere?**  } \\
\texttt{   "Nowhere" is the absence of place. It's not a location, but the concept of non{-}location. It could be emptiness, silence, or the void.} \\
            }
        }
    }
    & & \\ \\

    \theutterance \stepcounter{utterance}  
    & \multicolumn{4}{p{0.6\linewidth}}{
        \cellcolor[rgb]{0.9,0.9,0.9}{
            \makecell[{{p{\linewidth}}}]{
                \texttt{\tiny{[P1$\rangle$GM]}}
                \texttt{These questions are philosophical and playful, often used in riddles or meditative inquiry. Let's explore each one thoughtfully:} \\
\\ 
\texttt{1. **What is inside?**  } \\
\texttt{   This could refer to the inner self, consciousness, or the unseen aspects of things. In a physical sense, it might be the contents of an object. But in a deeper sense, "inside" can mean the essence or truth of something — what lies beneath the surface.  } \\
\texttt{   *Answer:* The truth, the self, or the unseen.} \\
\\ 
\texttt{2. **What is on your left?**  } \\
\texttt{   This is a literal question, but it depends on perspective. If you're facing forward, your left side is a direction. But metaphorically, it could be what you're not focusing on — the overlooked, the ignored, or the unseen.  } \\
\texttt{   *Answer:* The unseen, the overlooked, or simply "left" as in the direction.} \\
\\ 
\texttt{3. **What is here?**  } \\
\texttt{   "Here" refers to the present moment, the immediate space. It could be anything — air, light, silence, presence. But in a deeper sense, "here" is where everything is — the now.  } \\
\texttt{   *Answer:* The present moment, awareness, or existence.} \\
\\ 
\texttt{4. **What is nowhere?**  } \\
\texttt{   "Nowhere" is the absence of place. It's not a location, but the concept of non{-}location. It could be emptiness, silence, or the void.} \\
            }
        }
    }
    & & \\ \\

    \theutterance \stepcounter{utterance}  
    & & \multicolumn{4}{p{0.6\linewidth}}{
        \cellcolor[rgb]{0.9,0.9,0.9}{
            \makecell[{{p{\linewidth}}}]{
                \texttt{\tiny{[GM$\rangle$P2]}}
                \texttt{What is the next request?} \\
            }
        }
    }
    & \\ \\

    \theutterance \stepcounter{utterance}  
    & & & \multicolumn{4}{p{0.6\linewidth}}{
        \cellcolor[rgb]{0.9,0.9,0.9}{
            \makecell[{{p{\linewidth}}}]{
                \texttt{\tiny{[GM$\langle$P2]}}
                \texttt{What is on your right?} \\
            }
        }
    }
     \\ \\

    \theutterance \stepcounter{utterance}  
    & & \multicolumn{4}{p{0.6\linewidth}}{
        \cellcolor[rgb]{0.9,0.9,0.9}{
            \makecell[{{p{\linewidth}}}]{
                \texttt{\tiny{[P1$\langle$GM]}}
                \texttt{QUESTIONER: What is on your right?} \\
            }
        }
    }
    & \\ \\

    \theutterance \stepcounter{utterance}  
    & \multicolumn{4}{p{0.45\linewidth}}{
        \cellcolor[rgb]{0.95,0.95,0.95}{
            \makecell[{{p{\linewidth}}}]{
                \texttt{\tiny{[P1$|$internal]}}
                \texttt{{[}attempt 1 – valid{]}} \\
\texttt{These questions are playful and philosophical, often used in riddles or to provoke thought. Let's explore them one by one:} \\
\\ 
\texttt{1. **What is inside?**  } \\
\texttt{   — This depends on context. If you're asking about a box, a room, or a person, the answer varies. But in a riddle{-}like sense, "inside" could be *nothing* — like the inside of an empty container. Or it could be *the answer to the question itself* — a meta{-}answer.} \\
\\ 
\texttt{2. **What is on your left?**  } \\
\texttt{   — This is a physical question. If I were a person, it would be whatever is to my left. But since I'm an AI, I don't have a physical body. So, *nothing* — or perhaps *the question you just asked*.} \\
\\ 
\texttt{3. **What is here?**  } \\
\texttt{   — "Here" refers to the current place — in this case, the space of our conversation. So, *this chat*, *your questions*, *my responses*, and *the digital space* are all "here."} \\
\\ 
\texttt{4. **What is nowhere?**  } \\
\texttt{   — "Nowhere" is the absence of a place. So, *nothing* is "nowhere." Or, in a poetic sense, *the idea of non{-}existence* or *the void*.} \\
\\ 
\texttt{5. **What is there?**  } \\
\texttt{   — "There" refers to a place that is not here. In this} \\
            }
        }
    }
    & & \\ \\

    \theutterance \stepcounter{utterance}  
    & \multicolumn{4}{p{0.6\linewidth}}{
        \cellcolor[rgb]{0.9,0.9,0.9}{
            \makecell[{{p{\linewidth}}}]{
                \texttt{\tiny{[P1$\rangle$GM]}}
                \texttt{These questions are playful and philosophical, often used in riddles or to provoke thought. Let's explore them one by one:} \\
\\ 
\texttt{1. **What is inside?**  } \\
\texttt{   — This depends on context. If you're asking about a box, a room, or a person, the answer varies. But in a riddle{-}like sense, "inside" could be *nothing* — like the inside of an empty container. Or it could be *the answer to the question itself* — a meta{-}answer.} \\
\\ 
\texttt{2. **What is on your left?**  } \\
\texttt{   — This is a physical question. If I were a person, it would be whatever is to my left. But since I'm an AI, I don't have a physical body. So, *nothing* — or perhaps *the question you just asked*.} \\
\\ 
\texttt{3. **What is here?**  } \\
\texttt{   — "Here" refers to the current place — in this case, the space of our conversation. So, *this chat*, *your questions*, *my responses*, and *the digital space* are all "here."} \\
\\ 
\texttt{4. **What is nowhere?**  } \\
\texttt{   — "Nowhere" is the absence of a place. So, *nothing* is "nowhere." Or, in a poetic sense, *the idea of non{-}existence* or *the void*.} \\
\\ 
\texttt{5. **What is there?**  } \\
\texttt{   — "There" refers to a place that is not here. In this} \\
            }
        }
    }
    & & \\ \\

    \theutterance \stepcounter{utterance}  
    & & \multicolumn{4}{p{0.6\linewidth}}{
        \cellcolor[rgb]{0.9,0.9,0.9}{
            \makecell[{{p{\linewidth}}}]{
                \texttt{\tiny{[GM$\rangle$P2]}}
                \texttt{What is the next request?} \\
            }
        }
    }
    & \\ \\

    \theutterance \stepcounter{utterance}  
    & & & \multicolumn{4}{p{0.6\linewidth}}{
        \cellcolor[rgb]{0.9,0.9,0.9}{
            \makecell[{{p{\linewidth}}}]{
                \texttt{\tiny{[GM$\langle$P2]}}
                \texttt{What is in the southeast?} \\
            }
        }
    }
     \\ \\

    \theutterance \stepcounter{utterance}  
    & & \multicolumn{4}{p{0.6\linewidth}}{
        \cellcolor[rgb]{0.9,0.9,0.9}{
            \makecell[{{p{\linewidth}}}]{
                \texttt{\tiny{[P1$\langle$GM]}}
                \texttt{QUESTIONER: What is in the southeast?} \\
            }
        }
    }
    & \\ \\

    \theutterance \stepcounter{utterance}  
    & \multicolumn{4}{p{0.45\linewidth}}{
        \cellcolor[rgb]{0.95,0.95,0.95}{
            \makecell[{{p{\linewidth}}}]{
                \texttt{\tiny{[P1$|$internal]}}
                \texttt{{[}attempt 1 – valid{]}} \\
\texttt{These questions seem to be playful, philosophical, or possibly part of a riddle or game. Let's explore them one by one with a mix of literal, metaphorical, and logical interpretations:} \\
\\ 
\texttt{{-}{-}{-}} \\
\\ 
\texttt{**1. What is inside?**  } \\
\texttt{*Answer:* "Inside" refers to the space within a boundary. It could be anything—objects, feelings, thoughts, or even emptiness.  } \\
\texttt{**Possible answer:** "Whatever is enclosed by a boundary—like the air inside a balloon, or the thoughts inside a mind."} \\
\\ 
\texttt{{-}{-}{-}} \\
\\ 
\texttt{**2. What is on your left?**  } \\
\texttt{*Answer:* This depends on the context. If I'm facing forward, my left is to the side. But since I'm an AI, I don't have a physical body.  } \\
\texttt{**Possible answer:** "Nothing, because I don't have a left side. But metaphorically, it could be the unknown, the past, or the unseen."} \\
\\ 
\texttt{{-}{-}{-}} \\
\\ 
\texttt{**3. What is here?**  } \\
\texttt{*Answer:* "Here" refers to the current location—where the conversation is taking place.  } \\
\texttt{**Possible answer:** "The present moment, the dialogue, the space between us, and the digital environment where this exchange happens."} \\
\\ 
\texttt{{-}{-}{-}} \\
\\ 
\texttt{**4. What is nowhere?**  } \\
\texttt{*Answer:* "Nowhere" is the absence of a place. It's a concept, not a physical thing.  } \\
\texttt{**Possible answer:** "Nothingness, emptiness, or the idea of non{-}existence. It} \\
            }
        }
    }
    & & \\ \\

    \theutterance \stepcounter{utterance}  
    & \multicolumn{4}{p{0.6\linewidth}}{
        \cellcolor[rgb]{0.9,0.9,0.9}{
            \makecell[{{p{\linewidth}}}]{
                \texttt{\tiny{[P1$\rangle$GM]}}
                \texttt{These questions seem to be playful, philosophical, or possibly part of a riddle or game. Let's explore them one by one with a mix of literal, metaphorical, and logical interpretations:} \\
\\ 
\texttt{{-}{-}{-}} \\
\\ 
\texttt{**1. What is inside?**  } \\
\texttt{*Answer:* "Inside" refers to the space within a boundary. It could be anything—objects, feelings, thoughts, or even emptiness.  } \\
\texttt{**Possible answer:** "Whatever is enclosed by a boundary—like the air inside a balloon, or the thoughts inside a mind."} \\
\\ 
\texttt{{-}{-}{-}} \\
\\ 
\texttt{**2. What is on your left?**  } \\
\texttt{*Answer:* This depends on the context. If I'm facing forward, my left is to the side. But since I'm an AI, I don't have a physical body.  } \\
\texttt{**Possible answer:** "Nothing, because I don't have a left side. But metaphorically, it could be the unknown, the past, or the unseen."} \\
\\ 
\texttt{{-}{-}{-}} \\
\\ 
\texttt{**3. What is here?**  } \\
\texttt{*Answer:* "Here" refers to the current location—where the conversation is taking place.  } \\
\texttt{**Possible answer:** "The present moment, the dialogue, the space between us, and the digital environment where this exchange happens."} \\
\\ 
\texttt{{-}{-}{-}} \\
\\ 
\texttt{**4. What is nowhere?**  } \\
\texttt{*Answer:* "Nowhere" is the absence of a place. It's a concept, not a physical thing.  } \\
\texttt{**Possible answer:** "Nothingness, emptiness, or the idea of non{-}existence. It} \\
            }
        }
    }
    & & \\ \\

    \theutterance \stepcounter{utterance}  
    & & \multicolumn{4}{p{0.6\linewidth}}{
        \cellcolor[rgb]{0.9,0.9,0.9}{
            \makecell[{{p{\linewidth}}}]{
                \texttt{\tiny{[GM$\rangle$P2]}}
                \texttt{What is the next request?} \\
            }
        }
    }
    & \\ \\

    \theutterance \stepcounter{utterance}  
    & & & \multicolumn{4}{p{0.6\linewidth}}{
        \cellcolor[rgb]{0.9,0.9,0.9}{
            \makecell[{{p{\linewidth}}}]{
                \texttt{\tiny{[GM$\langle$P2]}}
                \texttt{What is in the southwest?} \\
            }
        }
    }
     \\ \\

    \theutterance \stepcounter{utterance}  
    & & \multicolumn{4}{p{0.6\linewidth}}{
        \cellcolor[rgb]{0.9,0.9,0.9}{
            \makecell[{{p{\linewidth}}}]{
                \texttt{\tiny{[P1$\langle$GM]}}
                \texttt{QUESTIONER: What is in the southwest?} \\
            }
        }
    }
    & \\ \\

    \theutterance \stepcounter{utterance}  
    & \multicolumn{4}{p{0.45\linewidth}}{
        \cellcolor[rgb]{0.95,0.95,0.95}{
            \makecell[{{p{\linewidth}}}]{
                \texttt{\tiny{[P1$|$internal]}}
                \texttt{{[}attempt 1 – valid{]}} \\
\texttt{These questions seem to be a playful or philosophical set, possibly inspired by riddles, meditation, or existential inquiry. Let's explore each one thoughtfully:} \\
\\ 
\texttt{1. **What is inside?**  } \\
\texttt{   — This could refer to the inner self, thoughts, emotions, or the unseen aspects of a person or object. In a physical sense, it might be the contents of a container. Metaphorically, it could be the soul, consciousness, or hidden truth.} \\
\\ 
\texttt{2. **What is on your left?**  } \\
\texttt{   — This depends on the context. If you're facing forward, it's the side to your left. In a symbolic sense, it might represent the past, the intuitive, or the feminine (in some traditions). In a literal sense, it could be a person, object, or direction.} \\
\\ 
\texttt{3. **What is here?**  } \\
\texttt{   — "Here" refers to the present moment and location. So, "what is here?" could be the immediate surroundings, the air, light, sound, or even the awareness of being present. It's the now.} \\
\\ 
\texttt{4. **What is nowhere?**  } \\
\texttt{   — "Nowhere" is the absence of place. It's not a location, so nothing is *in* nowhere. It's a concept of emptiness, absence, or nonexistence. It could also be a metaphor for the void, silence, or the unattainable.} \\
\\ 
\texttt{5. **What is there?**  } \\
\texttt{   —} \\
            }
        }
    }
    & & \\ \\

    \theutterance \stepcounter{utterance}  
    & \multicolumn{4}{p{0.6\linewidth}}{
        \cellcolor[rgb]{0.9,0.9,0.9}{
            \makecell[{{p{\linewidth}}}]{
                \texttt{\tiny{[P1$\rangle$GM]}}
                \texttt{These questions seem to be a playful or philosophical set, possibly inspired by riddles, meditation, or existential inquiry. Let's explore each one thoughtfully:} \\
\\ 
\texttt{1. **What is inside?**  } \\
\texttt{   — This could refer to the inner self, thoughts, emotions, or the unseen aspects of a person or object. In a physical sense, it might be the contents of a container. Metaphorically, it could be the soul, consciousness, or hidden truth.} \\
\\ 
\texttt{2. **What is on your left?**  } \\
\texttt{   — This depends on the context. If you're facing forward, it's the side to your left. In a symbolic sense, it might represent the past, the intuitive, or the feminine (in some traditions). In a literal sense, it could be a person, object, or direction.} \\
\\ 
\texttt{3. **What is here?**  } \\
\texttt{   — "Here" refers to the present moment and location. So, "what is here?" could be the immediate surroundings, the air, light, sound, or even the awareness of being present. It's the now.} \\
\\ 
\texttt{4. **What is nowhere?**  } \\
\texttt{   — "Nowhere" is the absence of place. It's not a location, so nothing is *in* nowhere. It's a concept of emptiness, absence, or nonexistence. It could also be a metaphor for the void, silence, or the unattainable.} \\
\\ 
\texttt{5. **What is there?**  } \\
\texttt{   —} \\
            }
        }
    }
    & & \\ \\

    \theutterance \stepcounter{utterance}  
    & & \multicolumn{4}{p{0.6\linewidth}}{
        \cellcolor[rgb]{0.9,0.9,0.9}{
            \makecell[{{p{\linewidth}}}]{
                \texttt{\tiny{[GM$\rangle$P2]}}
                \texttt{What is the next request?} \\
            }
        }
    }
    & \\ \\

    \theutterance \stepcounter{utterance}  
    & & & \multicolumn{4}{p{0.6\linewidth}}{
        \cellcolor[rgb]{0.9,0.9,0.9}{
            \makecell[{{p{\linewidth}}}]{
                \texttt{\tiny{[GM$\langle$P2]}}
                \texttt{What is outside?} \\
            }
        }
    }
     \\ \\

    \theutterance \stepcounter{utterance}  
    & & \multicolumn{4}{p{0.6\linewidth}}{
        \cellcolor[rgb]{0.9,0.9,0.9}{
            \makecell[{{p{\linewidth}}}]{
                \texttt{\tiny{[P1$\langle$GM]}}
                \texttt{QUESTIONER: What is outside?} \\
            }
        }
    }
    & \\ \\

    \theutterance \stepcounter{utterance}  
    & \multicolumn{4}{p{0.45\linewidth}}{
        \cellcolor[rgb]{0.95,0.95,0.95}{
            \makecell[{{p{\linewidth}}}]{
                \texttt{\tiny{[P1$|$internal]}}
                \texttt{{[}attempt 1 – valid{]}} \\
\texttt{These questions seem to be part of a riddle or a philosophical inquiry, possibly inspired by Zen koans or meditative puzzles. Let's explore them thoughtfully:} \\
\\ 
\texttt{{-}{-}{-}} \\
\\ 
\texttt{**1. What is inside?**  } \\
\texttt{*Answer:* The inside is what is enclosed, but also what is within us—thoughts, feelings, awareness. In a physical sense, it could be anything contained within a space. But in a deeper sense, "inside" may refer to the self, the mind, or the unseen.} \\
\\ 
\texttt{**2. What is on your left?**  } \\
\texttt{*Answer:* This depends on perspective. If you're facing forward, it's the left side. But metaphorically, it could be the unknown, the past, or the intuitive. In a riddle, it might be "nothing" or "the other side."} \\
\\ 
\texttt{**3. What is here?**  } \\
\texttt{*Answer:* "Here" is the present moment. It is where you are, now. It is the space of awareness. In Zen terms, "here" is everything and nothing—just what is.} \\
\\ 
\texttt{**4. What is nowhere?**  } \\
\texttt{*Answer:* "Nowhere" is the absence of place. It is not a location, but a concept. It could be emptiness, silence, or the void. In philosophy, it's the space beyond all places—perhaps the mind's stillness.} \\
\\ 
\texttt{**5. What is there?**  } \\
\texttt{*Answer:* "There" is the opposite of} \\
            }
        }
    }
    & & \\ \\

    \theutterance \stepcounter{utterance}  
    & \multicolumn{4}{p{0.6\linewidth}}{
        \cellcolor[rgb]{0.9,0.9,0.9}{
            \makecell[{{p{\linewidth}}}]{
                \texttt{\tiny{[P1$\rangle$GM]}}
                \texttt{These questions seem to be part of a riddle or a philosophical inquiry, possibly inspired by Zen koans or meditative puzzles. Let's explore them thoughtfully:} \\
\\ 
\texttt{{-}{-}{-}} \\
\\ 
\texttt{**1. What is inside?**  } \\
\texttt{*Answer:* The inside is what is enclosed, but also what is within us—thoughts, feelings, awareness. In a physical sense, it could be anything contained within a space. But in a deeper sense, "inside" may refer to the self, the mind, or the unseen.} \\
\\ 
\texttt{**2. What is on your left?**  } \\
\texttt{*Answer:* This depends on perspective. If you're facing forward, it's the left side. But metaphorically, it could be the unknown, the past, or the intuitive. In a riddle, it might be "nothing" or "the other side."} \\
\\ 
\texttt{**3. What is here?**  } \\
\texttt{*Answer:* "Here" is the present moment. It is where you are, now. It is the space of awareness. In Zen terms, "here" is everything and nothing—just what is.} \\
\\ 
\texttt{**4. What is nowhere?**  } \\
\texttt{*Answer:* "Nowhere" is the absence of place. It is not a location, but a concept. It could be emptiness, silence, or the void. In philosophy, it's the space beyond all places—perhaps the mind's stillness.} \\
\\ 
\texttt{**5. What is there?**  } \\
\texttt{*Answer:* "There" is the opposite of} \\
            }
        }
    }
    & & \\ \\

    \theutterance \stepcounter{utterance}  
    & & \multicolumn{4}{p{0.6\linewidth}}{
        \cellcolor[rgb]{0.9,0.9,0.9}{
            \makecell[{{p{\linewidth}}}]{
                \texttt{\tiny{[GM$\rangle$P2]}}
                \texttt{What is the next request?} \\
            }
        }
    }
    & \\ \\

    \theutterance \stepcounter{utterance}  
    & & & \multicolumn{4}{p{0.6\linewidth}}{
        \cellcolor[rgb]{0.9,0.9,0.9}{
            \makecell[{{p{\linewidth}}}]{
                \texttt{\tiny{[GM$\langle$P2]}}
                \texttt{What is in the northwest?} \\
            }
        }
    }
     \\ \\

    \theutterance \stepcounter{utterance}  
    & & \multicolumn{4}{p{0.6\linewidth}}{
        \cellcolor[rgb]{0.9,0.9,0.9}{
            \makecell[{{p{\linewidth}}}]{
                \texttt{\tiny{[P1$\langle$GM]}}
                \texttt{QUESTIONER: What is in the northwest?} \\
            }
        }
    }
    & \\ \\

    \theutterance \stepcounter{utterance}  
    & \multicolumn{4}{p{0.45\linewidth}}{
        \cellcolor[rgb]{0.95,0.95,0.95}{
            \makecell[{{p{\linewidth}}}]{
                \texttt{\tiny{[P1$|$internal]}}
                \texttt{{[}attempt 1 – valid{]}} \\
\texttt{These questions seem to be part of a riddle or a philosophical inquiry, possibly inspired by Zen koans or spatial awareness puzzles. Let's explore them thoughtfully:} \\
\\ 
\texttt{{-}{-}{-}} \\
\\ 
\texttt{**1. What is inside?**  } \\
\texttt{*Answer:* The answer depends on context. If you're asking about a room, a box, or a person, "inside" refers to what is contained within. But in a deeper sense, "inside" could mean the inner self, thoughts, or consciousness.  } \\
\texttt{*Possible answer:* **The mind** or **the self**.} \\
\\ 
\texttt{{-}{-}{-}} \\
\\ 
\texttt{**2. What is on your left?**  } \\
\texttt{*Answer:* This is a literal question about spatial orientation. If you're facing forward, it's the side to your left. But if this is a riddle, it might be pointing to something symbolic.  } \\
\texttt{*Possible answer:* **The unknown** or **the past** — or simply **your left hand**.} \\
\\ 
\texttt{{-}{-}{-}} \\
\\ 
\texttt{**3. What is here?**  } \\
\texttt{*Answer:* "Here" refers to the present location. In a physical sense, it could be the room, the air, the ground. In a metaphysical sense, it could be **presence**, **now**, or **awareness**.  } \\
\texttt{*Possible answer:* **Now** or **the present moment**.} \\
\\ 
\texttt{{-}{-}{-}} \\
\\ 
\texttt{**4. What is nowhere?**  } \\
\texttt{*Answer:* "Nowhere" is the absence of place. It's not a location, so nothing is there. But paradox} \\
            }
        }
    }
    & & \\ \\

    \theutterance \stepcounter{utterance}  
    & \multicolumn{4}{p{0.6\linewidth}}{
        \cellcolor[rgb]{0.9,0.9,0.9}{
            \makecell[{{p{\linewidth}}}]{
                \texttt{\tiny{[P1$\rangle$GM]}}
                \texttt{These questions seem to be part of a riddle or a philosophical inquiry, possibly inspired by Zen koans or spatial awareness puzzles. Let's explore them thoughtfully:} \\
\\ 
\texttt{{-}{-}{-}} \\
\\ 
\texttt{**1. What is inside?**  } \\
\texttt{*Answer:* The answer depends on context. If you're asking about a room, a box, or a person, "inside" refers to what is contained within. But in a deeper sense, "inside" could mean the inner self, thoughts, or consciousness.  } \\
\texttt{*Possible answer:* **The mind** or **the self**.} \\
\\ 
\texttt{{-}{-}{-}} \\
\\ 
\texttt{**2. What is on your left?**  } \\
\texttt{*Answer:* This is a literal question about spatial orientation. If you're facing forward, it's the side to your left. But if this is a riddle, it might be pointing to something symbolic.  } \\
\texttt{*Possible answer:* **The unknown** or **the past** — or simply **your left hand**.} \\
\\ 
\texttt{{-}{-}{-}} \\
\\ 
\texttt{**3. What is here?**  } \\
\texttt{*Answer:* "Here" refers to the present location. In a physical sense, it could be the room, the air, the ground. In a metaphysical sense, it could be **presence**, **now**, or **awareness**.  } \\
\texttt{*Possible answer:* **Now** or **the present moment**.} \\
\\ 
\texttt{{-}{-}{-}} \\
\\ 
\texttt{**4. What is nowhere?**  } \\
\texttt{*Answer:* "Nowhere" is the absence of place. It's not a location, so nothing is there. But paradox} \\
            }
        }
    }
    & & \\ \\

    \theutterance \stepcounter{utterance}  
    & & \multicolumn{4}{p{0.6\linewidth}}{
        \cellcolor[rgb]{0.9,0.9,0.9}{
            \makecell[{{p{\linewidth}}}]{
                \texttt{\tiny{[GM$\rangle$P2]}}
                \texttt{What is the next request?} \\
            }
        }
    }
    & \\ \\

    \theutterance \stepcounter{utterance}  
    & & & \multicolumn{4}{p{0.6\linewidth}}{
        \cellcolor[rgb]{0.9,0.9,0.9}{
            \makecell[{{p{\linewidth}}}]{
                \texttt{\tiny{[GM$\langle$P2]}}
                \texttt{What is at the center?} \\
            }
        }
    }
     \\ \\

    \theutterance \stepcounter{utterance}  
    & & \multicolumn{4}{p{0.6\linewidth}}{
        \cellcolor[rgb]{0.9,0.9,0.9}{
            \makecell[{{p{\linewidth}}}]{
                \texttt{\tiny{[P1$\langle$GM]}}
                \texttt{QUESTIONER: What is at the center?} \\
            }
        }
    }
    & \\ \\

    \theutterance \stepcounter{utterance}  
    & \multicolumn{4}{p{0.45\linewidth}}{
        \cellcolor[rgb]{0.95,0.95,0.95}{
            \makecell[{{p{\linewidth}}}]{
                \texttt{\tiny{[P1$|$internal]}}
                \texttt{{[}attempt 1 – valid{]}} \\
\texttt{These questions seem to be part of a riddle or a philosophical inquiry, possibly inspired by Zen koans or spatial awareness puzzles. Let's explore them thoughtfully:} \\
\\ 
\texttt{{-}{-}{-}} \\
\\ 
\texttt{**1. What is inside?**  } \\
\texttt{→ *The answer depends on context.*  } \\
\texttt{If you're asking about a room, a box, or a mind — it could be anything. But in a metaphorical sense, "inside" might refer to the self, thoughts, or hidden truths.  } \\
\texttt{**Answer:** *The unknown, the inner world, or whatever is enclosed.*} \\
\\ 
\texttt{{-}{-}{-}} \\
\\ 
\texttt{**2. What is on your left?**  } \\
\texttt{→ *This is a physical question.*  } \\
\texttt{If I were a person, I’d look to my left. But as an AI, I don’t have a physical body.  } \\
\texttt{**Answer:** *Nothing — or the space where a left side would be, if I had one.*} \\
\\ 
\texttt{{-}{-}{-}} \\
\\ 
\texttt{**3. What is here?**  } \\
\texttt{→ *This refers to the present moment, the location of the question.*  } \\
\texttt{**Answer:** *The space where the question is being asked — the now, the present, the interaction.*} \\
\\ 
\texttt{{-}{-}{-}} \\
\\ 
\texttt{**4. What is nowhere?**  } \\
\texttt{→ *A paradoxical question.*  } \\
\texttt{"Where" is a place, but "nowhere" is the absence of place.  } \\
\texttt{**Answer:** *Nothing — or the concept of non{-}existence, emptiness, or the void.*} \\
\\ 
\texttt{{-}{-}{-}} \\
\\ 
\texttt{**5. What is there?**  } \\
\texttt{→ *Opposite} \\
            }
        }
    }
    & & \\ \\

    \theutterance \stepcounter{utterance}  
    & \multicolumn{4}{p{0.6\linewidth}}{
        \cellcolor[rgb]{0.9,0.9,0.9}{
            \makecell[{{p{\linewidth}}}]{
                \texttt{\tiny{[P1$\rangle$GM]}}
                \texttt{These questions seem to be part of a riddle or a philosophical inquiry, possibly inspired by Zen koans or spatial awareness puzzles. Let's explore them thoughtfully:} \\
\\ 
\texttt{{-}{-}{-}} \\
\\ 
\texttt{**1. What is inside?**  } \\
\texttt{→ *The answer depends on context.*  } \\
\texttt{If you're asking about a room, a box, or a mind — it could be anything. But in a metaphorical sense, "inside" might refer to the self, thoughts, or hidden truths.  } \\
\texttt{**Answer:** *The unknown, the inner world, or whatever is enclosed.*} \\
\\ 
\texttt{{-}{-}{-}} \\
\\ 
\texttt{**2. What is on your left?**  } \\
\texttt{→ *This is a physical question.*  } \\
\texttt{If I were a person, I’d look to my left. But as an AI, I don’t have a physical body.  } \\
\texttt{**Answer:** *Nothing — or the space where a left side would be, if I had one.*} \\
\\ 
\texttt{{-}{-}{-}} \\
\\ 
\texttt{**3. What is here?**  } \\
\texttt{→ *This refers to the present moment, the location of the question.*  } \\
\texttt{**Answer:** *The space where the question is being asked — the now, the present, the interaction.*} \\
\\ 
\texttt{{-}{-}{-}} \\
\\ 
\texttt{**4. What is nowhere?**  } \\
\texttt{→ *A paradoxical question.*  } \\
\texttt{"Where" is a place, but "nowhere" is the absence of place.  } \\
\texttt{**Answer:** *Nothing — or the concept of non{-}existence, emptiness, or the void.*} \\
\\ 
\texttt{{-}{-}{-}} \\
\\ 
\texttt{**5. What is there?**  } \\
\texttt{→ *Opposite} \\
            }
        }
    }
    & & \\ \\

    \theutterance \stepcounter{utterance}  
    & & \multicolumn{4}{p{0.6\linewidth}}{
        \cellcolor[rgb]{0.9,0.9,0.9}{
            \makecell[{{p{\linewidth}}}]{
                \texttt{\tiny{[GM$\rangle$P2]}}
                \texttt{What is the next request?} \\
            }
        }
    }
    & \\ \\

    \theutterance \stepcounter{utterance}  
    & & & \multicolumn{4}{p{0.6\linewidth}}{
        \cellcolor[rgb]{0.9,0.9,0.9}{
            \makecell[{{p{\linewidth}}}]{
                \texttt{\tiny{[GM$\langle$P2]}}
                \texttt{What is in the northeast?} \\
            }
        }
    }
     \\ \\

    \theutterance \stepcounter{utterance}  
    & & \multicolumn{4}{p{0.6\linewidth}}{
        \cellcolor[rgb]{0.9,0.9,0.9}{
            \makecell[{{p{\linewidth}}}]{
                \texttt{\tiny{[P1$\langle$GM]}}
                \texttt{QUESTIONER: What is in the northeast?} \\
            }
        }
    }
    & \\ \\

    \theutterance \stepcounter{utterance}  
    & \multicolumn{4}{p{0.45\linewidth}}{
        \cellcolor[rgb]{0.95,0.95,0.95}{
            \makecell[{{p{\linewidth}}}]{
                \texttt{\tiny{[P1$|$internal]}}
                \texttt{{[}attempt 1 – valid{]}} \\
\texttt{These questions seem to be part of a riddle or a philosophical/meditative exercise, possibly inspired by traditional koans or spatial awareness puzzles. Let's explore them thoughtfully.} \\
\\ 
\texttt{{-}{-}{-}} \\
\\ 
\texttt{**1. What is inside?**  } \\
\texttt{→ *The answer depends on context.*  } \\
\texttt{If you're inside a room, it's the space enclosed by walls.  } \\
\texttt{If you're inside a person, it's organs, thoughts, or consciousness.  } \\
\texttt{But in a riddle{-}like sense: **"Inside"** could be **"nothing"** — because everything is *outside* of the self, or **"the mind"** — as the inner world.} \\
\\ 
\texttt{**2. What is on your left?**  } \\
\texttt{→ *Depends on your orientation.*  } \\
\texttt{But in a symbolic sense: **"The past"** — because left is often associated with memory or the past.} \\
\\ 
\texttt{**3. What is here?**  } \\
\texttt{→ *The present moment.*  } \\
\texttt{"Here" refers to the current location — but also to **"now"** — the only real time we experience.} \\
\\ 
\texttt{**4. What is nowhere?**  } \\
\texttt{→ *The answer is "nowhere" itself.*  } \\
\texttt{It’s a paradox: **"Nowhere" is the absence of place.**  } \\
\texttt{But in Zen terms: **"Nowhere" is everywhere** — because emptiness is the ground of all being.} \\
\\ 
\texttt{**5. What is there?**  } \\
\texttt{→ *The opposite of "here" — the unknown} \\
            }
        }
    }
    & & \\ \\

    \theutterance \stepcounter{utterance}  
    & \multicolumn{4}{p{0.6\linewidth}}{
        \cellcolor[rgb]{0.9,0.9,0.9}{
            \makecell[{{p{\linewidth}}}]{
                \texttt{\tiny{[P1$\rangle$GM]}}
                \texttt{These questions seem to be part of a riddle or a philosophical/meditative exercise, possibly inspired by traditional koans or spatial awareness puzzles. Let's explore them thoughtfully.} \\
\\ 
\texttt{{-}{-}{-}} \\
\\ 
\texttt{**1. What is inside?**  } \\
\texttt{→ *The answer depends on context.*  } \\
\texttt{If you're inside a room, it's the space enclosed by walls.  } \\
\texttt{If you're inside a person, it's organs, thoughts, or consciousness.  } \\
\texttt{But in a riddle{-}like sense: **"Inside"** could be **"nothing"** — because everything is *outside* of the self, or **"the mind"** — as the inner world.} \\
\\ 
\texttt{**2. What is on your left?**  } \\
\texttt{→ *Depends on your orientation.*  } \\
\texttt{But in a symbolic sense: **"The past"** — because left is often associated with memory or the past.} \\
\\ 
\texttt{**3. What is here?**  } \\
\texttt{→ *The present moment.*  } \\
\texttt{"Here" refers to the current location — but also to **"now"** — the only real time we experience.} \\
\\ 
\texttt{**4. What is nowhere?**  } \\
\texttt{→ *The answer is "nowhere" itself.*  } \\
\texttt{It’s a paradox: **"Nowhere" is the absence of place.**  } \\
\texttt{But in Zen terms: **"Nowhere" is everywhere** — because emptiness is the ground of all being.} \\
\\ 
\texttt{**5. What is there?**  } \\
\texttt{→ *The opposite of "here" — the unknown} \\
            }
        }
    }
    & & \\ \\

    \theutterance \stepcounter{utterance}  
    & & \multicolumn{4}{p{0.6\linewidth}}{
        \cellcolor[rgb]{0.9,0.9,0.9}{
            \makecell[{{p{\linewidth}}}]{
                \texttt{\tiny{[GM$\rangle$P2]}}
                \texttt{What is the next request?} \\
            }
        }
    }
    & \\ \\

    \theutterance \stepcounter{utterance}  
    & & & \multicolumn{4}{p{0.6\linewidth}}{
        \cellcolor[rgb]{0.9,0.9,0.9}{
            \makecell[{{p{\linewidth}}}]{
                \texttt{\tiny{[GM$\langle$P2]}}
                \texttt{What is at the bottom?} \\
            }
        }
    }
     \\ \\

    \theutterance \stepcounter{utterance}  
    & & \multicolumn{4}{p{0.6\linewidth}}{
        \cellcolor[rgb]{0.9,0.9,0.9}{
            \makecell[{{p{\linewidth}}}]{
                \texttt{\tiny{[P1$\langle$GM]}}
                \texttt{QUESTIONER: What is at the bottom?} \\
            }
        }
    }
    & \\ \\

    \theutterance \stepcounter{utterance}  
    & \multicolumn{4}{p{0.45\linewidth}}{
        \cellcolor[rgb]{0.95,0.95,0.95}{
            \makecell[{{p{\linewidth}}}]{
                \texttt{\tiny{[P1$|$internal]}}
                \texttt{{[}attempt 1 – valid{]}} \\
\texttt{These questions seem to be part of a riddle or a philosophical/meditative exercise, possibly inspired by spatial awareness, perception, or Zen koans. Let's explore them thoughtfully:} \\
\\ 
\texttt{{-}{-}{-}} \\
\\ 
\texttt{**1. What is inside?**  } \\
\texttt{→ *The answer depends on context.*  } \\
\texttt{If you're inside a room: the air, objects, people, thoughts.  } \\
\texttt{If you're inside a body: organs, blood, consciousness.  } \\
\texttt{But in a deeper sense: *everything that is not outside* — perhaps *the self*, or *awareness*.} \\
\\ 
\texttt{**2. What is on your left?**  } \\
\texttt{→ *Depends on your orientation.*  } \\
\texttt{But metaphorically: *the unknown*, *the past*, *the intuitive*, *the feminine* (in some traditions).  } \\
\texttt{Or simply: *what you haven't yet turned to*.} \\
\\ 
\texttt{**3. What is here?**  } \\
\texttt{→ *Now*, *this moment*, *presence*.  } \\
\texttt{The answer is often: *everything that is present* — the air, light, sound, your breath, the question itself.} \\
\\ 
\texttt{**4. What is nowhere?**  } \\
\texttt{→ *Nothing*, *void*, *emptiness*, *the absence of place*.  } \\
\texttt{But also: *the space between thoughts*, *the ungraspable*, *the infinite*.} \\
\\ 
\texttt{**5. What is there?**  } \\
\texttt{→ *The opposite of here* — *the distant*, *the unseen*, *the future*, *the unknown*.} \\
            }
        }
    }
    & & \\ \\

    \theutterance \stepcounter{utterance}  
    & \multicolumn{4}{p{0.6\linewidth}}{
        \cellcolor[rgb]{0.9,0.9,0.9}{
            \makecell[{{p{\linewidth}}}]{
                \texttt{\tiny{[P1$\rangle$GM]}}
                \texttt{These questions seem to be part of a riddle or a philosophical/meditative exercise, possibly inspired by spatial awareness, perception, or Zen koans. Let's explore them thoughtfully:} \\
\\ 
\texttt{{-}{-}{-}} \\
\\ 
\texttt{**1. What is inside?**  } \\
\texttt{→ *The answer depends on context.*  } \\
\texttt{If you're inside a room: the air, objects, people, thoughts.  } \\
\texttt{If you're inside a body: organs, blood, consciousness.  } \\
\texttt{But in a deeper sense: *everything that is not outside* — perhaps *the self*, or *awareness*.} \\
\\ 
\texttt{**2. What is on your left?**  } \\
\texttt{→ *Depends on your orientation.*  } \\
\texttt{But metaphorically: *the unknown*, *the past*, *the intuitive*, *the feminine* (in some traditions).  } \\
\texttt{Or simply: *what you haven't yet turned to*.} \\
\\ 
\texttt{**3. What is here?**  } \\
\texttt{→ *Now*, *this moment*, *presence*.  } \\
\texttt{The answer is often: *everything that is present* — the air, light, sound, your breath, the question itself.} \\
\\ 
\texttt{**4. What is nowhere?**  } \\
\texttt{→ *Nothing*, *void*, *emptiness*, *the absence of place*.  } \\
\texttt{But also: *the space between thoughts*, *the ungraspable*, *the infinite*.} \\
\\ 
\texttt{**5. What is there?**  } \\
\texttt{→ *The opposite of here* — *the distant*, *the unseen*, *the future*, *the unknown*.} \\
            }
        }
    }
    & & \\ \\

    \theutterance \stepcounter{utterance}  
    & & \multicolumn{4}{p{0.6\linewidth}}{
        \cellcolor[rgb]{0.9,0.9,0.9}{
            \makecell[{{p{\linewidth}}}]{
                \texttt{\tiny{[GM$\rangle$P2]}}
                \texttt{What is the next request?} \\
            }
        }
    }
    & \\ \\

    \theutterance \stepcounter{utterance}  
    & & & \multicolumn{4}{p{0.6\linewidth}}{
        \cellcolor[rgb]{0.9,0.9,0.9}{
            \makecell[{{p{\linewidth}}}]{
                \texttt{\tiny{[GM$\langle$P2]}}
                \texttt{What is everywhere?} \\
            }
        }
    }
     \\ \\

    \theutterance \stepcounter{utterance}  
    & & \multicolumn{4}{p{0.6\linewidth}}{
        \cellcolor[rgb]{0.9,0.9,0.9}{
            \makecell[{{p{\linewidth}}}]{
                \texttt{\tiny{[P1$\langle$GM]}}
                \texttt{QUESTIONER: What is everywhere?} \\
            }
        }
    }
    & \\ \\

    \theutterance \stepcounter{utterance}  
    & \multicolumn{4}{p{0.45\linewidth}}{
        \cellcolor[rgb]{0.95,0.95,0.95}{
            \makecell[{{p{\linewidth}}}]{
                \texttt{\tiny{[P1$|$internal]}}
                \texttt{{[}attempt 1 – valid{]}} \\
\texttt{These questions are poetic and philosophical, playing with the concepts of location, presence, absence, and perception. Here's a thoughtful interpretation of each:} \\
\\ 
\texttt{{-}{-}{-}} \\
\\ 
\texttt{**1. What is inside?**  } \\
\texttt{*Answer:* The unseen, the hidden, the core — what is contained within boundaries, both physical and metaphorical. It could be the heart, the mind, or the essence of something.} \\
\\ 
\texttt{**2. What is on your left?**  } \\
\texttt{*Answer:* The unknown, the past, the shadow — what lies to the side, not directly in front, often overlooked or unacknowledged.} \\
\\ 
\texttt{**3. What is here?**  } \\
\texttt{*Answer:* The present moment — the now, the immediate reality, the space where you are standing.} \\
\\ 
\texttt{**4. What is nowhere?**  } \\
\texttt{*Answer:* Absence, emptiness, the void — the concept of non{-}existence, or the place where nothing is.} \\
\\ 
\texttt{**5. What is there?**  } \\
\texttt{*Answer:* The distant, the future, the unseen — what exists beyond the immediate, perhaps in the horizon or in thought.} \\
\\ 
\texttt{**6. What is on your right?**  } \\
\texttt{*Answer:* The future, the unknown, the side not chosen — what lies in the path not taken, or the opposite of the left.} \\
\\ 
\texttt{**7. What is in the southeast?**  } \\
\texttt{*Answer:* The warmth, the light, the convergence of elements — a direction often associated with growth, energy, and the sun.} \\
\\ 
\texttt{**8} \\
            }
        }
    }
    & & \\ \\

    \theutterance \stepcounter{utterance}  
    & \multicolumn{4}{p{0.6\linewidth}}{
        \cellcolor[rgb]{0.9,0.9,0.9}{
            \makecell[{{p{\linewidth}}}]{
                \texttt{\tiny{[P1$\rangle$GM]}}
                \texttt{These questions are poetic and philosophical, playing with the concepts of location, presence, absence, and perception. Here's a thoughtful interpretation of each:} \\
\\ 
\texttt{{-}{-}{-}} \\
\\ 
\texttt{**1. What is inside?**  } \\
\texttt{*Answer:* The unseen, the hidden, the core — what is contained within boundaries, both physical and metaphorical. It could be the heart, the mind, or the essence of something.} \\
\\ 
\texttt{**2. What is on your left?**  } \\
\texttt{*Answer:* The unknown, the past, the shadow — what lies to the side, not directly in front, often overlooked or unacknowledged.} \\
\\ 
\texttt{**3. What is here?**  } \\
\texttt{*Answer:* The present moment — the now, the immediate reality, the space where you are standing.} \\
\\ 
\texttt{**4. What is nowhere?**  } \\
\texttt{*Answer:* Absence, emptiness, the void — the concept of non{-}existence, or the place where nothing is.} \\
\\ 
\texttt{**5. What is there?**  } \\
\texttt{*Answer:* The distant, the future, the unseen — what exists beyond the immediate, perhaps in the horizon or in thought.} \\
\\ 
\texttt{**6. What is on your right?**  } \\
\texttt{*Answer:* The future, the unknown, the side not chosen — what lies in the path not taken, or the opposite of the left.} \\
\\ 
\texttt{**7. What is in the southeast?**  } \\
\texttt{*Answer:* The warmth, the light, the convergence of elements — a direction often associated with growth, energy, and the sun.} \\
\\ 
\texttt{**8} \\
            }
        }
    }
    & & \\ \\

    \theutterance \stepcounter{utterance}  
    & & \multicolumn{4}{p{0.6\linewidth}}{
        \cellcolor[rgb]{0.9,0.9,0.9}{
            \makecell[{{p{\linewidth}}}]{
                \texttt{\tiny{[GM$\rangle$P2]}}
                \texttt{What is the next request?} \\
            }
        }
    }
    & \\ \\

    \theutterance \stepcounter{utterance}  
    & & & \multicolumn{4}{p{0.6\linewidth}}{
        \cellcolor[rgb]{0.9,0.9,0.9}{
            \makecell[{{p{\linewidth}}}]{
                \texttt{\tiny{[GM$\langle$P2]}}
                \texttt{What is at the top?} \\
            }
        }
    }
     \\ \\

    \theutterance \stepcounter{utterance}  
    & & \multicolumn{4}{p{0.6\linewidth}}{
        \cellcolor[rgb]{0.9,0.9,0.9}{
            \makecell[{{p{\linewidth}}}]{
                \texttt{\tiny{[P1$\langle$GM]}}
                \texttt{QUESTIONER: What is at the top?} \\
            }
        }
    }
    & \\ \\

    \theutterance \stepcounter{utterance}  
    & \multicolumn{4}{p{0.45\linewidth}}{
        \cellcolor[rgb]{0.95,0.95,0.95}{
            \makecell[{{p{\linewidth}}}]{
                \texttt{\tiny{[P1$|$internal]}}
                \texttt{{[}attempt 1 – valid{]}} \\
\texttt{These questions seem to be a playful or philosophical exploration of spatial and existential concepts—possibly inspired by riddles, meditation, or a game like "Where's Waldo?" or a Zen koan. Let's answer them in a way that blends literal, metaphorical, and philosophical interpretations.} \\
\\ 
\texttt{{-}{-}{-}} \\
\\ 
\texttt{**1. What is inside?**  } \\
\texttt{*Inside* refers to what is contained within a boundary.  } \\
\texttt{→ **Answer:** The space enclosed by a container, structure, or system.  } \\
\texttt{*Philosophical twist:* The self, consciousness, or the unknown.} \\
\\ 
\texttt{{-}{-}{-}} \\
\\ 
\texttt{**2. What is on your left?**  } \\
\texttt{This depends on the observer’s orientation.  } \\
\texttt{→ **Answer:** The side of the body opposite the right.  } \\
\texttt{*Metaphor:* The past, the unknown, the subconscious.} \\
\\ 
\texttt{{-}{-}{-}} \\
\\ 
\texttt{**3. What is here?**  } \\
\texttt{"Here" is the current location of the speaker.  } \\
\texttt{→ **Answer:** The present moment, the immediate surroundings.  } \\
\texttt{*Philosophical twist:* Reality, presence, the now.} \\
\\ 
\texttt{{-}{-}{-}} \\
\\ 
\texttt{**4. What is nowhere?**  } \\
\texttt{"Nowhere" is the absence of a place.  } \\
\texttt{→ **Answer:** Nothing, emptiness, nonexistence.  } \\
\texttt{*Metaphor:* The void, silence, the unmanifest.} \\
\\ 
\texttt{{-}{-}{-}} \\
\\ 
\texttt{**5. What is there?**  } \\
\texttt{"There" is a location opposite "here."  } \\
\texttt{→ **Answer:** The place where something exists, but not where the speaker is.  } \\
\texttt{*Philosoph} \\
            }
        }
    }
    & & \\ \\

    \theutterance \stepcounter{utterance}  
    & \multicolumn{4}{p{0.6\linewidth}}{
        \cellcolor[rgb]{0.9,0.9,0.9}{
            \makecell[{{p{\linewidth}}}]{
                \texttt{\tiny{[P1$\rangle$GM]}}
                \texttt{These questions seem to be a playful or philosophical exploration of spatial and existential concepts—possibly inspired by riddles, meditation, or a game like "Where's Waldo?" or a Zen koan. Let's answer them in a way that blends literal, metaphorical, and philosophical interpretations.} \\
\\ 
\texttt{{-}{-}{-}} \\
\\ 
\texttt{**1. What is inside?**  } \\
\texttt{*Inside* refers to what is contained within a boundary.  } \\
\texttt{→ **Answer:** The space enclosed by a container, structure, or system.  } \\
\texttt{*Philosophical twist:* The self, consciousness, or the unknown.} \\
\\ 
\texttt{{-}{-}{-}} \\
\\ 
\texttt{**2. What is on your left?**  } \\
\texttt{This depends on the observer’s orientation.  } \\
\texttt{→ **Answer:** The side of the body opposite the right.  } \\
\texttt{*Metaphor:* The past, the unknown, the subconscious.} \\
\\ 
\texttt{{-}{-}{-}} \\
\\ 
\texttt{**3. What is here?**  } \\
\texttt{"Here" is the current location of the speaker.  } \\
\texttt{→ **Answer:** The present moment, the immediate surroundings.  } \\
\texttt{*Philosophical twist:* Reality, presence, the now.} \\
\\ 
\texttt{{-}{-}{-}} \\
\\ 
\texttt{**4. What is nowhere?**  } \\
\texttt{"Nowhere" is the absence of a place.  } \\
\texttt{→ **Answer:** Nothing, emptiness, nonexistence.  } \\
\texttt{*Metaphor:* The void, silence, the unmanifest.} \\
\\ 
\texttt{{-}{-}{-}} \\
\\ 
\texttt{**5. What is there?**  } \\
\texttt{"There" is a location opposite "here."  } \\
\texttt{→ **Answer:** The place where something exists, but not where the speaker is.  } \\
\texttt{*Philosoph} \\
            }
        }
    }
    & & \\ \\

    \theutterance \stepcounter{utterance}  
    & & & \multicolumn{2}{p{0.3\linewidth}}{
        \cellcolor[rgb]{0.9,0.9,0.9}{
            \makecell[{{p{\linewidth}}}]{
                \texttt{\tiny{[GM$|$GM]}}
                \texttt{Game finished.} \\
            }
        }
    }
    & & \\ \\

\end{supertabular}
}

\end{document}
